\documentclass[../main.tex]{subfiles}

\begin{document}

\subsection{How the Evidence Builds the Argument}
\label{subsec:storytelling-principles}

The argument is built through accumulation rather than assertion. Each page
begins with a high-level anchor, then moves into a more diagnostic chart that
either qualifies or complicates the headline. For example, Page 2 does not stop
at high internet-use rates of \(87.85\%\) and \(85.39\%\); it contrasts them
with Active Digital ID Verification at only \(32.68\%\), then deepens the point
through Figure~\ref{fig:page2-digital-activity-profile}, where Social Media
reaches \(94.46\%\) but Gov activity is only \(15.63\%\). That sequence turns a
connectivity story into an activation story.

The same evidence-building logic continues in later pages. Page 3 moves from
the \(74.93\%\) banking headline to social segmentation and exclusion barriers.
Page 4 moves from digital wage inflows at \(80.99\%\) to the much weaker rural
and agricultural usage environment. Page 5 moves from a national Saved rate of
\(73.05\%\) to the more revealing contrast between Family \& Friends at 86 and
Formal Savings at 30 for the Poorest 20\%. Because each visual layer narrows
the meaning of the previous one, the reader is led toward a more cautious and
policy-useful interpretation of inclusion.

\noindent\textit{See Figures~\ref{fig:page2-readiness-overview},
\ref{fig:page2-digital-activity-profile}, \ref{fig:account-kpi-cards},
\ref{fig:page4-utility-payment-methods}, and
\ref{fig:page5-coping-income-group}.}

\end{document}
